\subsection{The contribution of the stratification of culture-talk for status-based inequality}  
\label{sec:inequality}

Having defined and established the status of culture talk as a theoretically relevant phenomenon, in this section I propose three substantive reasons for why culture talk is relevant for our understanding of the process the linkage between culture and stratification.  The gist of the argument is that culture talk is an important micro-mechanism in generating often observed empirical associations between status group inequalities premised on access to cultural and material resources and outcomes associated with the large-scale structuration of social ties and group formation in the contemporary context.  

Shedding empirical light on the way in which cultural tastes are input into substantively important network outcomes has been a relatively neglected issue in the literature on the culture/stratification link.  However, the connection between lifestyle consumption behaviors and network connectivity has begun to receive sustained empirical and conceptual attention of late \citep{erickson96, lizardo06a, wimmer_lewis10}.   Here I link this recent literature to classical theoretical statements and long-standing empirical generalizations connecting social capital and network outcomes to stratification and status attainment processes.  This research is consistent with the empirical expectation that inequalities in culture talk should be predictable (e.g. stratified by the familiar markers of status), pervasive, and consequential in a practical and explanatory sense.

\subsubsection{Relationship formation}
\label{subsubsec:relform}

First, culture talk allows for the formation and renewal of social relationships.  This is a ``need'' that is structurally induced by the fact that, as recent research on the dynamics of personal networks show, new relationships are constantly being formed and old ones renewed (or dropped) during the adult life-cycle \citep{wellman_etal97, suitor_keeton97, suitor_etal97, morgan_etal97, degenne_dynamics_2005, bidart_lavenu05, bidart_degenne05}.  The most surprising thing about ``intimacy'' and ``friendship'' is thus, how unstable it tends to be \citep{wellman_etal97}.  Cultural knowledge associated with aesthetic consumption experiences can thus be thought of as a key resource with which to form \emph{new} social relationships.  Culture talk is especially useful at the \emph{initial} stages of ``the acquaintance process'' \citep{newcomb61} allowing the potential contacts to search for deeper areas of commonality and agreement without much risk (as would inquiring about politics or economic position), so that the interaction can proceed smoothly, thus facilitating the initiation and continuation of conversational rituals when more intimate content is not yet available (and even after this content is available).  When individuals do not have a previous social or cultural basis for interaction based on common membership in the same social network or common biographical experiences they instead rely ``\ldots upon common cultural topics on which to build a conversation'' \citet[460]{fine77}. 

\subsubsection{Status-based solidarity in small groups}
\label{subsec:solidarity}

Second, culture talk serves is a key mechanism in the generation and reproduction of the micro-solidarities that hold status-groups together \citep{ostrower98}.  The common knowledge gained through culture talk serves as a cognitive and emotional basis of solidarity within the micro-ritual generated in the interaction, as persons become entranced by ``the joy of sharing names'' \citep{dimaggio87}. culture talk``\ldots facilitates the sociable intercourse necessary for acquaintanceships to ripen into friendships'' \citep{dimaggio87}.  Following this line of reasoning, we can begin to appreciate the ``social'' implications of the consumption of symbolic goods. These implications are intimately related to the maintenance and enhancement of status-linked social cohesion in small groups \citep{fine77}.  Classical theoretical statements in the culture and stratification literature argue that status-based exclusion occurs largely through micro-mechanisms governing entry (or determining systematic exclusion from) class-structured micro-worlds organized around specific conversational interactions \citep{collins75, collins00}.  The micro-solidarities that form the basis of status-linked orders but also form of the core of what has come to be called ``social capital'' are may thus be built around the most seemingly innocuous forms of conversational subject matter.  These theorists argue that the experiential and phenomenological ``stuff'' of stratification is essentially based on talk, and much of this talk is about culture. As John Fiske's observed, civic membership and ``belonging'' to neighborhoods and communities is based on talk, and ``[m]uch of this talk is about the mass media and its cultural commodities'' \citeyearpar[73]{fiske87}.  This is the reason why we can consider culture-talk to be a primary  ``means of construction social relations\ldots.''  Culture-talk helps persons to ``establish networks of trusting relations that facilitate group mobilization and the attainment of such social rewards as desirable spouses and prestigious jobs''  \citep[443]{dimaggio87}.  Culture-talk serves to increase the cohesion and social closure of small groups, providing a way to segment them from the larger social structure;  this is why we can think of culture-talk as being much more ``\ldots than just an acquaintanceship device.''  It is also ``a mechanism for establishing, confirming and reifying group beliefs and attitudes'' \citep[461]{fine77}.  Accordingly, one of the primary behavioral implications of routine access to aesthetic consumption experiences is precisely their recruitment as the non-instrumental content that routinely re-creates the existence of status-segregated conversational rituals:  ``\ldots talk about entertainment-the discussions of movies, music, sports, actors and athletes\ldots make up so much of casual chitchat.'' This becomes even more important ``\ldots if we include talk about things people consume or seek out-food, clothes, houses and gardens, restaurants scenery, vacations.  In some social circles, this entertainment talk makes up most of the content of conversations" \citep[125]{collins75}. For instance, Simmelian urban social circles are produced and reproduced through shared \emph{cultural tastes}--e.g. concert going, visiting the museum, etc.--and other activities intrinsically connected to culture-talk \citep[797]{kadushin66}.

\subsubsection{Linkages across social positions}
\label{subsubsec:linkages}
Third, certain forms of culture talk represent a key mediator in allowing for sporadic and sometimes sustained connections across socially distant positions in the social structure.  This is what has been referred in a recent paper by \citet{schultz_breiger10} as ``the strength of weak culture.''  One of the most significant transformations in the social logic that governs status in contemporary (post)industrial societies consists of the rising role of the arts, and mass produced culture in general, in providing the ``baseline'' forms of cultural knowledge necessary to sustain interaction for persons who cannot come to expect commonalities based on traditional symbols of acquaintance and recognition. This shift has occurred in the context of the social transformation of status-based collectivities from solidary groups with a strong allegiance to community and place, to looser, wide-spanning ``personal networks'' characterized by increasing geographic and social mobility \citep{wellman_etal88}.  This process acquires more importance as these network ties have been transformed in the contemporary context of increasing geographic mobility \citep{griswold_wright04} and the decline of the traditional bonds characteristic of traditional local communities \citep{wellman79}, which used to be centered around kinship and spatial contiguity \citep{dimaggio_mohr85}.  As the personal networks of socially mobile members of the upper-middle class have become more fluid, loosely knit, more discontinuous and less tied to spatial contiguity and kinship \citep{wellman_wortley90, wellman_etal88}, the arts and popular culture have come to increasingly serve as the ``default" forms of knowledge that connect people across different interaction foci.  culture talk that makes use of knowledge associated with symbolic goods that are truly popular may cut across class and status lines may in this way serve to connect persons across different positions in the social structure and across different taste publics, facilitating coordination across hierarchical positions \citep{erickson96}.   culture talk allows persons to create and sustain social connections--even if sometimes brief and sporadic (so-called ``weak ties'')--when there exist large socio-demographic differences between them \citep{lizardo06a}.

